%%%%%%%%%%%%%%%%%%%%%%%%%%%%%%%%%%%%%%%%%%%%%%%%%%%%%%%%%%%%
%
% Authors:       Nina Wiedemann, Benjamin Ummenhofer, Diana Wofk, 
%                Quentin Leboutet, and Michael Paulitsch
%
% Description:   The SYCL of Life: Evolutionary Kernel Optimization 
%                with Large Language Models
%
% Section:       Related Works
%
%%%%%%%%%%%%%%%%%%%%%%%%%%%%%%%%%%%%%%%%%%%%%%%%%%%%%%%%%%%%

\section{Related Work} \label{sec:related}
\paragraph{Traditional Kernel Optimization.}
High-performance computing has long relied on separating algorithm specification from optimization schedules, as exemplified in  Halide~\citep{ragan2013halide} or TVM~\citep{chen2018tvm}, whereas lower-level models like CUDA and SYCL require both to be handled together in code. Triton~\citep{tillet2019triton} offers a middle ground with a Python-based DSL that simplifies custom GPU kernel development but still exposes scheduling control.
Vendor libraries (cuDNN~\citep{chetlur2014cudnn}, CUTLASS~\citep{cutlass}, oneDNN~\citep{onednn}) deliver peak performance but lack flexibility for novel operators. Classical autotuners like ATLAS~\citep{whaley2001automated} and FFTW~\citep{frigo1998fftw} demonstrate that systematic search can match hand-tuned code. To avoid exhaustively evaluating every schedule, systems like AutoTVM~\citep{chen2018learning}, AnsorAnsor~\citep{zheng2020ansor}, and Ithemal~\citep{mendis2019ithemal} use learned cost models to predict performance, but these require extensive training data and struggle to generalize across hardware. Despite these tools, algorithmic kernel development remained a highly manual and expertise-driven process.

\paragraph{LLM-Based Kernel Generation.}
Large language models have sparked intense interest in automated kernel synthesis, yet generating \textit{correct} GPU code remains challenging, and generating \textit{fast} code is harder still. Early studies~\citep{valero2023comparing, godoy2023evaluation} evaluated Llama-2 and Codex on generating HPC kernels. The introduction of KernelBench~\citep{ouyang2025kernelbench} established the de facto benchmark, comprising 250~tasks spanning single operators, fusion patterns, and complete architectures. Results showed that even frontier models achieve correctness on merely ${\sim}20\%$ of problems with substantial gaps to cuBLAS/cuDNN baselines. Recent work has also highlighted validation challenges: spurious speedups of up to 50x can arise from incorrect kernels that pass test cases~\citep{lange2025towards, zhu2025establishing}, motivating more rigorous evaluation protocols. 
Despite these challenges, KernelBench has inspired a plethora of follow-up approaches (see Table~\ref{tab:related_work} for a comprehensive overview). 
The dominant paradigm is \textit{prompting with validation loop}: correctness and profiling outputs guide successive refinements of 
CUDA~\citep{chen2025cuda, zhang2025cudaforge, lange2025towards}, Triton~\citep{liao2025kernelevolve, li2025tritonforge}, and Metal~\citep{gimlet} kernels. PEAK~\citep{tariq2025peak} expresses optimization strategies as natural language transformations, achieving up to 95\% of cuBLAS performance on MatMul across CUDA, HIP, and HLSL. Multi-agent architectures decompose the task into specialized roles: Astra~\citep{wei2025astra} separates planning, coding, testing, and profiling agents; STARK~\citep{dong2025stark} uses plan-code-debug agents; QiMeng~\citep{zhu2025qimeng} employs macro/micro coding agents; and PRAGMA~\citep{lei2025pragma} coordinates coder, verifier, and conductor roles. SwizzlePerf~\citep{tschand2025swizzleperf} stands out by injecting detailed microarchitecture knowledge (bank conflicts, warp operations, coalescing) into prompts, demonstrating the value of hardware awareness. Despite their sophistication, these approaches share a common limitation: they explore the kernel design space implicitly through generic code generation, without mechanisms to preserve diversity across optimization trajectories.

An orthogonal direction is \textit{finetuning} LLMs for kernel generation. CUDA-L1~\citep{li2025cuda} combines supervised finetuning (SFT) with ``contrastive RL''; CUDA-L2~\citep{su2025cuda} adds NCU profiler metrics and multi-stage RL, exceeding cuBLASLt by 11\% on HGEMM, TritonRL~\citep{woo2025tritonrl} presents good results with SFT and GRPO for Triton, and Kevin~\citep{baronio2025kevin} uses multi-turn RL to refine kernels across dialogue turns. These methods achieve promising results but often cannot match the performance of the latest closed-source LLMs.

Existing LLM kernel work targets CUDA almost exclusively. MultiKernelBench~\citep{wen2025multikernelbench} addresses this with benchmarks spanning NVIDIA GPUs, Huawei NPUs, and Google TPUs, and % revealing poor generalization to under-represented platforms, and 
NPUEval~\citep{kalade2025npueval} provides AMD NPU benchmarks. Translation systems like BabelTower~\citep{wen2022babeltower} and CodeRosetta~\citep{tehrani2024coderosetta} convert between CUDA and HIP/OpenCL but preserve rather than optimize performance. Our work contributes to this direction by targeting SYCL~\citep{keryell2015khronos} and hence enabling portability across Intel, NVIDIA, and AMD accelerators. 

\paragraph{LLM-Guided Evolutionary Search.}
Evolutionary algorithms have a rich history in program synthesis~\citep{geneticprogramming}. Quality-diversity (QD) methods~\citep{pugh2016quality} shift from single-objective optimization to maintaining \textit{diverse collections} of high-performing solutions: a property particularly relevant for kernel optimization, where multiple valid implementation strategies exist. MAP-Elites~\citep{mouret2015illuminating} partitions the solution space by behavioral descriptors, CMA-ME~\citep{cmame} adds continuous optimization, and PGA-MAP-Elites~\citep{pgamapelites} integrates policy gradients. Recent work combines LLMs with evolutionary search: FunSearch~\citep{romera2024mathematical} discovers mathematical programs beyond human baselines, AlphaEvolve~\citep{novikov2025alphaevolve} adds multi-objective optimization, and CodeEvolve~\citep{codeevolve} evolves prompts rather than code directly.
PACEvolve~\citep{yan2026pacevolve} characterizes three failure modes in LLM-assisted evolution: \textit{context pollution} from accumulated failures, \textit{mode collapse} from poor exploration-exploitation balance, and \textit{weak collaboration} in multi-island settings. While PACEvolve addresses these through hierarchical context management for general search, kernel optimization presents a \textit{structured} space where domain-specific dimensions can be directly exploited. Our approach leverages this structure: kernel-specific behavioral coordinates (memory patterns, parallelism strategies) inherently partition the solution space, preventing mode collapse by construction, while meta-prompt evolution~\citep{gepa} counteracts context degradation by enabling prompts to specialize for different optimization regions.

